\documentclass[10pt,a4paper]{amsart}
\usepackage{amsmath,amssymb,amsthm,graphicx}
% Retained preamble for the Stacks Project chapter snapshot in this directory.
% The chapter's own preamble.tex is a copy of the shared Stacks Project preamble; the block
% below is sliced MECHANICALLY, line for line, from it. Nothing was added, removed or
% redefined; the wrapper class, the xr-hyper/hyperref wiring and the externaldocument list
% are omitted because the assembler supplies the document class.
%
% lines taken from preamble.tex: 3, 5, 9, 11, 12, 156, 157, 158, ... (51 lines total)
%
% package coverage restored from the chapter's own preamble (2026-10-13): xy, xy-2cell

\usepackage[all]{xy}
\xyoption{2cell}
\UseAllTwocells

\newenvironment{reference}{\comment}{\endcomment}
\newenvironment{slogan}{\comment}{\endcomment}
\newenvironment{history}{\comment}{\endcomment}
\theoremstyle{plain}
\newtheorem{theorem}[subsection]{Theorem}
\newtheorem{proposition}[subsection]{Proposition}
\newtheorem{lemma}[subsection]{Lemma}
\theoremstyle{definition}
\newtheorem{definition}[subsection]{Definition}
\newtheorem{example}[subsection]{Example}
\newtheorem{exercise}[subsection]{Exercise}
\newtheorem{situation}[subsection]{Situation}
\theoremstyle{remark}
\newtheorem{remark}[subsection]{Remark}
\newtheorem{remarks}[subsection]{Remarks}
\numberwithin{equation}{subsection}
\def\lim{\mathop{\mathrm{lim}}\nolimits}
\def\colim{\mathop{\mathrm{colim}}\nolimits}
\def\Spec{\mathop{\mathrm{Spec}}}
\def\Hom{\mathop{\mathrm{Hom}}\nolimits}
\def\Ext{\mathop{\mathrm{Ext}}\nolimits}
\def\SheafHom{\mathop{\mathcal{H}\!\mathit{om}}\nolimits}
\def\SheafExt{\mathop{\mathcal{E}\!\mathit{xt}}\nolimits}
\def\Sch{\mathit{Sch}}
\def\Mor{\mathop{\mathrm{Mor}}\nolimits}
\def\Ob{\mathop{\mathrm{Ob}}\nolimits}
\def\Sh{\mathop{\mathit{Sh}}\nolimits}
\def\NL{\mathop{N\!L}\nolimits}
\def\CH{\mathop{\mathrm{CH}}\nolimits}
\def\proetale{{pro\text{-}\acute{e}tale}}
\def\etale{{\acute{e}tale}}
\def\QCoh{\mathit{QCoh}}
\def\Ker{\mathop{\mathrm{Ker}}}
\def\Im{\mathop{\mathrm{Im}}}
\def\Coker{\mathop{\mathrm{Coker}}}
\def\Coim{\mathop{\mathrm{Coim}}}
\DeclareMathSymbol{\boxtimes}{\mathbin}{AMSa}{"02}
\def\QCohstack{\mathcal{QC}\!\mathit{oh}}
\def\Cohstack{\mathcal{C}\!\mathit{oh}}
\def\Spacesstack{\mathcal{S}\!\mathit{paces}}
\def\Quotfunctor{\mathrm{Quot}}
\def\Hilbfunctor{\mathrm{Hilb}}
\def\Curvesstack{\mathcal{C}\!\mathit{urves}}
\def\Polarizedstack{\mathcal{P}\!\mathit{olarized}}
\def\Complexesstack{\mathcal{C}\!\mathit{omplexes}}
\def\Pic{\mathop{\mathrm{Pic}}\nolimits}
\def\Picardstack{\mathcal{P}\!\mathit{ic}}
\def\Picardfunctor{\mathrm{Pic}}
\def\Deformationcategory{\mathcal{D}\!\mathit{ef}}
\usepackage{cleveref}
\crefname{originallemma}{original lemma}{original lemmas}
\Crefname{originallemma}{Original lemma}{Original lemmas}
\makeatletter
\newlabel{algebraic-lemma-open-fibred-category-is-algebraic}{{05UN}{source}{}{sourcepointer-1}{}}
\newlabel{algebraic-lemma-open-fibred-category-is-algebraic@cref}{{[originallemma][0][]05UN}{[1][0][]source}{}{}{}}
\makeatother
\crefname{originallemma}{original lemma}{original lemmas}
\Crefname{originallemma}{Original lemma}{Original lemmas}
\makeatletter
\newlabel{stacks-lemma-2-fibre-product-gives-stack-in-setoids}{{05UI}{source}{}{sourcepointer-2}{}}
\newlabel{stacks-lemma-2-fibre-product-gives-stack-in-setoids@cref}{{[originallemma][0][]05UI}{[1][0][]source}{}{}{}}
\makeatother
\title{An open immersion of a presentation lifts to an open substack receiving it surjectively smoothly}
\author{The Stacks Project Authors}
\date{stacks-project:stacks-properties.tex:local-snapshot-20261012; GNU FDL 1.2 or later}
\begin{document}
\maketitle
\noindent\textit{Editorial note.} Statement span is lines 2084-2091 of sources/stacks-stacks-properties/pinned-original-source.tex; proof span is lines 2093-2143, the complete author proof (49 body lines).\par\medskip
\noindent\textit{Editorial note.} Navigation stubs: 2 cited-but-not-included labels are supplied as non-load-bearing pointers with their pinned tags where the snapshot tag index has them.\par\medskip
\noindent\textit{Editorial note.} extra\_packages is empty on purpose: the retained preamble already loads xy with options and a second declaration clashes in the assembler wrapper.\par\medskip
\noindent\textit{Editorial note.} Text is retained verbatim, including the author's own typos and placeholders.\par\medskip
\noindent\textit{Editorial note.} Queue index 2 of queue\_stacks-properties.json (pb=2271, stmt 2084-2091, proof 2093-2143).\par\medskip
\noindent\textit{Editorial note.} stacks-properties wave-1: fresh adjudication accept (2026-10-14); pool scanned while CONTROL.md is down (0 bytes) - §2 row deferred per Chief, claim reported to Chief 08:0x.\par\medskip
\noindent\textit{Allowed standard prerequisites (editorial).}\begin{itemize}
\item Explicit prerequisite boundary. The cited inputs are prerequisites; the argument in the retained span is not replaced by any cited result.
\end{itemize}\medskip
\ifdefined\hypertarget\hypertarget{sourcepointer-1}{}\fi
\noindent\textit{Source reference.} Original lemma 05UN: sources/stacks-stacks-properties/pinned-original-source.tex, source lines 2143--2143. Cited by the retained statement or proof but not part of the collected spans; supplied only as a navigation pointer. The argument does not depend on its content.\par\medskip
\ifdefined\hypertarget\hypertarget{sourcepointer-2}{}\fi
\noindent\textit{Source reference.} Original lemma 05UI: sources/stacks-stacks-properties/pinned-original-source.tex, source lines 2143--2143. Cited by the retained statement or proof but not part of the collected spans; supplied only as a navigation pointer. The argument does not depend on its content.\par\medskip
\begin{lemma}
\label{lemma-open-image-substack}
Let $\mathcal X$ be an algebraic stack. Let $U$ be an algebraic space and
$U \to \mathcal X$ a surjective smooth morphism. For an open immersion
$V \hookrightarrow U$, there exists an algebraic stack $\mathcal Y$, an
open immersion $\mathcal Y \to \mathcal X$, and a surjective smooth
morphism $V \to \mathcal Y$.
\end{lemma}

\begin{proof}
We define a category fibred in groupoids $\mathcal Y$ by letting the fiber
category $\mathcal{Y}_T$ over an object $T$ of $(\Sch/S)_{fppf}$ be
the full subcategory of $\mathcal{X}_T$ consisting of all
$y \in \Ob(\mathcal{X}_T)$ such that the projection morphism
$V \times_{\mathcal X, y} T \to T$ surjective. Now for any morphism
$x : T \to \mathcal X$, the $2$-fibred product
$T \times_{x, \mathcal X} \mathcal Y$ has fiber category over $T'$ consisting
of triples $(f : T' \to T, y \in \mathcal{X}_{T'}, f^*x \simeq y)$ such that
$V \times_{\mathcal X, y} T' \to T'$ is surjective.
Note that $T \times_{x, \mathcal X} \mathcal Y$ is fibered in setoids since
$\mathcal Y \to \mathcal X$ is faithful (see
Stacks, Lemma \ref{stacks-lemma-2-fibre-product-gives-stack-in-setoids}).
Now the isomorphism $f^*x \simeq y$ gives the diagram
$$
\xymatrix{
V \times_{\mathcal X, y} T' \ar[d] \ar[r] &
V \times_{\mathcal X, x} T \ar[r] \ar[d] &
V \ar[d] \\
T' \ar[r]^f &
T \ar[r]^x &
\mathcal X
}
$$
where both squares are cartesian. The morphism
$V \times_{\mathcal X, x} T \to T$ is smooth by base change, and hence open.
Let $T_0 \subset T$ be its image. From the cartesian squares we deduce that
$V \times_{\mathcal X, y} T' \to T'$ is surjective if and only if $f$ lands
in $T_0$.  Therefore $T \times_{x, \mathcal X} \mathcal Y$ is representable by
$T_0$, so the inclusion $\mathcal Y \to \mathcal X$ is an open immersion.
By
Algebraic Stacks, Lemma \ref{algebraic-lemma-open-fibred-category-is-algebraic}
we conclude that $\mathcal{Y}$ is an algebraic stack.
Lastly if we denote the morphism $V \to \mathcal X$ by $g$, we have
$V \times_{\mathcal X} V \to V$ is surjective (the diagonal gives a
section). Hence $g$ is in the image of $\mathcal{Y}_V \to \mathcal{X}_V$, i.e.,
we obtain a morphism $g' : V \to \mathcal{Y}$ fitting into the commutative
diagram
$$
\xymatrix{
V \ar[r] \ar[d]^{g'} & U \ar[d] \\
\mathcal{Y} \ar[r] & \mathcal{X}
}
$$
Since $V \times_{g, \mathcal X} \mathcal Y \to V$ is a
monomorphism, it is in fact an isomorphism since $(1, g')$ defines a section.
Therefore $g' : V \to \mathcal Y$ is a smooth morphism, as it is the
base change of the smooth morphism $g : V \to \mathcal{X}$.
It is surjective by our construction of $\mathcal{Y}$ which finishes
the proof of the lemma.
\end{proof}

\end{document}
